\documentclass[11pt]{article}
\usepackage[margin=1in]{geometry}
\usepackage{amsmath,amssymb,amsthm}
\usepackage{graphicx}
\usepackage{microtype}
\usepackage[hidelinks]{hyperref}
\usepackage{enumitem}

\newtheorem{theorem}{Theorem}
\newtheorem{lemma}{Lemma}
\newcommand{\R}{\mathbb R}
\newcommand{\N}{\mathbb N}
\newcommand{\1}{\mathbf 1}
\newcommand{\eps}{\varepsilon}

\title{An incomplete abstract Hardy space\\
\large Counterexample to the unqualified completeness question in arXiv:0809.4110}
\author{Automatic Mathematics Run \texttt{fa\_banach\_001}, lane 11}
\date{9 August 2026}

\begin{document}
\maketitle

\begin{center}
\fbox{\parbox{0.91\textwidth}{\textbf{Status.} Candidate counterexample, likely valid, pending human review. The construction gives a genuine normed but incomplete atomic/molecular Hardy space under the source paper's stated minimal hypotheses.}}
\end{center}

\section{Source question}

The source is F.~Bernicot, \emph{Use of abstract Hardy spaces, Real interpolation and Applications to bilinear operators}, arXiv:0809.4110v1 \cite{Bernicot}. In Remark~2.6(1), page~6, the author writes that the molecular expansion is only required to be defined almost everywhere and asks whether the resulting abstract normed spaces are complete.

\begin{figure}[ht]
  \centering
  \includegraphics[width=\textwidth]{figures/open_problem_crop.png}
  \caption{Remark~2.6 on page~6 of the source paper.}
\end{figure}

The exact sentence is: ``It is not clear whether these abstract normed vector spaces are complete.'' The construction below answers this negatively under the hypotheses stated before the remark.

\section{Definitions used from the source}

Let $(X,d,\mu)$ be a space of homogeneous type, let $\beta\in(1,\infty]$, and let $(B_Q)_Q$ be linear operators indexed by balls and uniformly bounded on $L^\beta(X)$. For $\beta=2$, an atom associated with a ball $Q$ is $B_Qf$, where $f$ is supported in $Q$ and
\[
  \|f\|_2\leq \mu(Q)^{-1/2}.
\]
An $\eps$-molecule has the same form with the source paper's annular decay bounds; on the innermost annulus these again give
\[
  \|f\|_{L^2(Q)}\leq \mu(Q)^{-1/2}.
\]
The Hardy gauge of a measurable $h$ is the infimum of $\sum_j|\lambda_j|$ over almost-everywhere convergent representations $h=\sum_j\lambda_jm_j$ by atoms or molecules.

\section{Proof intuition}

The only assumption on the family $(B_Q)$ is uniform $L^2$ boundedness. A very small ball therefore permits a normalized input of very large $L^2$ norm. We use a rank-one $B_{Q_n}$ to turn that allowance into the large atom $n^2r_n$, where $r_n$ is the $n$th Rademacher function.

An atomic expansion then has two simultaneous requirements. Its coefficient cost is the weak weighted quantity $\sum |b_n|/n^2$, while almost-everywhere convergence of the Rademacher series forces the strong condition $b\in\ell_2$. Thus the Hardy space carries only the weighted $\ell_1$ norm on the proper dense subspace
\[
  \ell_2\cap\ell_1(n^{-2})\subsetneq\ell_1(n^{-2}).
\]
The truncations of $(1,1,\ldots)$ are Cauchy for the displayed norm but cannot converge inside the space.

\section{Counterexample}

\begin{theorem}\label{thm:main}
For every $\eps\in(0,\infty]$ there exist a space of homogeneous type $(X,d,\mu)$, an exponent $\beta=2$, and a uniformly $L^2$-bounded family $(B_Q)_Q$ such that the source paper's spaces satisfy
\[
  H^1_{\eps,\mathrm{mol}}=H^1_{\mathrm{ato}}
\]
isometrically and are normed but incomplete.

More precisely, the common space is isometric to
\[
  E:=\left\{b=(b_n)\in\ell_2:\ \sum_{n=1}^\infty\frac{|b_n|}{n^2}<\infty\right\},
  \qquad
  \|b\|_E:=\sum_{n=1}^\infty\frac{|b_n|}{n^2}.
\]
\end{theorem}

\begin{proof}
Take $X=(0,1)$ with Euclidean distance and Lebesgue measure. This is a space of homogeneous type. Let $(r_n)_{n\geq1}$ be the standard Rademacher functions on $(0,1)$, so $|r_n|=1$ almost everywhere and the $r_n$ are independent and orthonormal in $L^2(0,1)$.

For $n\geq1$, let
\[
  Q_n=\left(0,\frac1{n^4}\right)
\]
(viewed as the open ball with center $1/(2n^4)$ and radius $1/(2n^4)$), and set
\[
  e_n:=\mu(Q_n)^{-1/2}\1_{Q_n}=n^2\1_{Q_n}.
\]
Define
\[
  B_{Q_n}f:=\langle f,e_n\rangle_{L^2}r_n,
  \qquad B_Q:=0\quad\text{for every other ball }Q.
\]
Since $\|e_n\|_2=\|r_n\|_2=1$, every nonzero $B_{Q_n}$ has operator norm one. Hence $(B_Q)_Q$ satisfies the required uniform $L^2$ bound.

Put $f_n=n^4\1_{Q_n}=n^2e_n$. Then $f_n$ is supported in $Q_n$ and
\[
  \|f_n\|_2=n^2=\mu(Q_n)^{-1/2},
  \qquad B_{Q_n}f_n=n^2r_n=:a_n.
\]
Thus $a_n$ is an atom. Conversely, if $m=B_{Q_n}f$ is either an atom or an $\eps$-molecule, its innermost-annulus bound gives
\[
  |\langle f,e_n\rangle|
  \leq \|f\|_{L^2(Q_n)}\|e_n\|_2
  \leq n^2.
\]
Because $B_{Q_n}$ only reads $f$ on $Q_n$, the outer molecular bounds do not enlarge the possible outputs. The nonzero normalized atoms and $\eps$-molecules are therefore exactly the functions $t a_n$ with $|t|\leq1$. In particular, all the atomic and molecular spaces and gauges coincide.

We now identify this common space. Consider an admissible representation
\[
  h=\sum_{j=1}^\infty \alpha_j a_{n_j}\qquad\text{almost everywhere},
  \qquad \sum_j|\alpha_j|<\infty,
\]
where smaller normalized atoms have been absorbed into $\alpha_j$. Let
\[
  S_N:=\sum_{j=1}^N\alpha_jn_j^2r_{n_j}
       =\sum_{n\geq1}c_n^{(N)}r_n,
\]
where $c^{(N)}$ has finite support. Since $S_N$ converges almost everywhere, it is Cauchy in probability.

For every finite scalar sequence $u$, the Rademacher sum $Z=\sum_nu_nr_n$ satisfies
\[
  \mathbb E Z^2=\|u\|_2^2,
  \qquad \mathbb E Z^4\leq3\|u\|_2^4.
\]
Applying Paley--Zygmund to $Z^2$ gives
\[
  \mathbb P\left(|Z|\geq\frac{\|u\|_2}{\sqrt2}\right)\geq\frac1{12}.
\]
Consequently, Cauchy convergence in probability of the finite sums $S_N$ forces $(c^{(N)})$ to be Cauchy in $\ell_2$. Let its limit be $b\in\ell_2$. Coordinatewise convergence and $\sum_j|\alpha_j|<\infty$ give
\[
  b_n=n^2\sum_{j:n_j=n}\alpha_j,
  \qquad
  \sum_n\frac{|b_n|}{n^2}\leq\sum_j|\alpha_j|.
\]
Moreover $S_N\to\sum_nb_nr_n$ in $L^2$, and this $L^2$ limit agrees almost everywhere with $h$.

Conversely, if $b\in E$, then
\[
  \sum_{n=1}^\infty\frac{b_n}{n^2}a_n=\sum_{n=1}^\infty b_nr_n
\]
has summable atomic coefficients. Its partial sums form an $L^2$-bounded martingale and hence converge almost everywhere (equivalently one may use the standard convergence theorem for square-summable Rademacher series; see \cite[Chapter 22]{Billingsley}). Therefore this is an admissible representation. The preceding lower bound and this representation show
\[
  \left\|\sum_nb_nr_n\right\|_{H^1}
  =\sum_n\frac{|b_n|}{n^2}.
\]
This also proves that the gauge is a genuine norm: a zero function has all of its unique $L^2$ Rademacher coefficients equal to zero.

Finally define
\[
  h_N:=\sum_{n=1}^N r_n.
\]
For $M>N$ the norm identity gives
\[
  \|h_M-h_N\|_{H^1}=\sum_{n=N+1}^M\frac1{n^2}\longrightarrow0,
\]
so $(h_N)$ is Cauchy. If it converged in $H^1$ to the element represented by $b\in E$, then
\[
  \sum_{n=1}^\infty\frac{|b_n-\1_{\{n\leq N\}}|}{n^2}\longrightarrow0.
\]
For each fixed $n$, taking $N\geq n$ and then $N\to\infty$ yields $b_n=1$. This contradicts $b\in\ell_2$. Hence the Hardy space is incomplete.
\end{proof}

\section{Verification and scope}

The finite checker in \texttt{code/finite\_rademacher\_check.py} verifies exact finite-dyadic orthogonality, the fourth-moment inequality for all vectors in $\{-1,0,1\}^6$, and decay of the weighted Cauchy tails. It is a sanity check only; the proof above is analytic.

The counterexample uses a deliberately sparse, discontinuous-in-$Q$ family of rank-one operators. It settles the universal question under the stated assumptions, but it does not say that familiar Hardy spaces arising from semigroups, approximations of the identity, or geometrically compatible cancellation operators are incomplete. It also does not conflict with Proposition~2.8 of the source, whose additional continuous embedding into $L^1_{\mathrm{loc}}$ rules out this example.

\section{Novelty check}

A bounded web/arXiv search on 9 August 2026 used the exact sentence from Remark~2.6, the source title and arXiv id, the author/title combination with ``complete atoms molecules'', and the notation $H^1_{\eps,\mathrm{mol}}$. The search returned the source paper and related abstract Hardy-space papers, but no later explicit resolution and no equivalent Rademacher/rank-one construction. This supports, but does not prove, novelty; confidence is moderate.

\section{Human-review recommendation}

The main review points are: (i) the source framework truly imposes no compatibility in $Q$ beyond uniform $L^2$ boundedness; (ii) the selected intervals are balls in the subspace $(0,1)$; (iii) the molecular output set equals the atomic output set; and (iv) the Paley--Zygmund argument handles arbitrary repetition and ordering of atoms. Subject to those checks, this should be accepted as a full negative answer at the source's stated level of generality.

\begin{thebibliography}{9}
\bibitem{Bernicot}
F.~Bernicot,
\emph{Use of abstract Hardy spaces, Real interpolation and Applications to bilinear operators},
arXiv:0809.4110v1 (2008).

\bibitem{Billingsley}
P.~Billingsley,
\emph{Probability and Measure}, 3rd ed.,
Wiley, 1995.
\end{thebibliography}

\end{document}
